\par\Needspace{6\baselineskip}\medskip
\begin{center}\hypertarget{native-178-article131}{}\pdfbookmark[2]{131.}{article-131}131.\end{center}
Article 114 showed that if the ratios among the quantities denoted there and in Article 128 by
$n,n',n''$ are known, the celestial body's distances from Earth can be determined by very simple formulas.
If, then, we took for $x,y$ the quotients $\dfrac{n}{n'},\ \dfrac{n''}{n'}$, and the heliocentric motion between observations were
not very large, the approximate values $\dfrac{\theta}{\theta'},\ \dfrac{\theta''}{\theta'}$ would immediately suggest themselves, with the symbols
$\theta,\theta',\theta''$ understood as in Article 128. A straightforward solution to our problem would thus seem to
follow by deriving two distances from Earth from $x$ and $y$, and then proceeding according to one
of the five methods in Articles 124–128. Indeed, if we likewise take the symbols $\eta,\eta''$ in the sense
of Article 128 and by analogy denote by $\eta'$ the quotient of the area of the sector enclosed by two
radius vectors and the area of the triangle between them, we obtain
\[
\frac{n}{n'}=\frac{\theta}{\theta'}\cdot\frac{\eta'}{\eta}\mathpunct{\raisebox{.45ex}{,}}\qquad
\frac{n''}{n'}=\frac{\theta''}{\theta'}\cdot\frac{\eta'}{\eta''}\mathpunct{\raisebox{.45ex}{,}}
\]
and if $n,n',n''$ are regarded as small quantities of the first order, it is easy to see that $\eta-1,\eta'-1,\eta''-1$ are
generally of the second order, and hence that the approximate values for $x,y$, namely $\dfrac{\theta}{\theta'},\ \dfrac{\theta''}{\theta'}$, differ
from the true ones only by quantities of the second order. On closer examination, however, this
method proves entirely unsuitable, for reasons we shall explain briefly. It is easy to see that the
quantity $(0.1.2)$, by which the distances are multiplied in formulas 9, 10, and 11 of Article 114, is of
at least the third order. In equation 9, for example, the quantities $(\mathrm{O}.1.2)$, $(\mathrm{I}.1.2)$, and $(\mathrm{II}.1.2)$ are of the
first order. It follows readily that a second-order error in the values of $\dfrac{n}{n'},\ \dfrac{n''}{n'}$ produces an error of
order $0$ in the values of the distances. Thus, in ordinary language, the resulting distances would
be affected by a finite error even when the time intervals are infinitely small. Neither these
distances nor the remaining quantities to be derived from them could be regarded even as
approximations, and the method would fail to satisfy the second condition of the preceding article.
